%-----------------------------------------------%
\usepackage[a4paper, top=2.5cm, bottom=2.5cm, left=2.5cm, right=2.5cm]{geometry}
\usepackage{amsmath,esint,mathtools}
\usepackage{indentfirst}
\usepackage{enumitem}
\usepackage{adjustbox}
\usepackage{float}
\usepackage{hyperref}
\usepackage{xcolor}
\usepackage{graphicx}
\usepackage{caption}
\usepackage{subcaption}
\usepackage{xltxtra}
\usepackage{multirow}
\usepackage{array}
\usepackage{titlesec}   % for \titleformat, \titlespacing
\usepackage{etoolbox}   % for \pretocmd
%------------ สำหรับการทำ Gantt Chart -----------%
\usepackage{colortbl}
\usepackage{forloop}
\newcounter{loopcntr}
\newcommand{\rpt}[2][1]{%
  \forloop{loopcntr}{0}{\value{loopcntr}<#1}{#2}%
}
\newcommand{\on}[1][1]{
  \forloop{loopcntr}{0}{\value{loopcntr}<#1}{&\cellcolor{lightgray}}
}
\newcommand{\off}[1][1]{
  \forloop{loopcntr}{0}{\value{loopcntr}<#1}{&}
}



\newlength{\firstindentlength}
\newlength{\secondindentlength}
\newlength{\thirdindentlength}

\setlength{\firstindentlength}{3em}
\setlength{\secondindentlength}{7.1em}
\setlength{\thirdindentlength}{12.2em}

%------------ สำหรับการใช้ฟ้อนท์ภาษาไทย -----------%
\usepackage[no-math]{fontspec} % สำหรับเลือก font
\usepackage{polyglossia}     % สำหรับหลายภาษา
\usepackage{xunicode}

\XeTeXlinebreaklocale "th"
\XeTeXlinebreakskip = 0pt plus 1pt

\renewcommand{\baselinestretch}{1.6}
\setmainfont[Script=Thai,Scale=1.4,ItalicFont={THSarabunNew_Italic.ttf}, BoldFont={THSarabunNew_Bold.ttf},BoldItalicFont={THSarabunNew_BoldItalic.ttf} ]{THSarabunNew.ttf}
\setmonofont[Script=Thai,Scale=0.9]{DroidSansMono.ttf}
%--------------- เปลี่ยนชื่อรูปและตาราง -------------%
\usepackage[labelsep=space]{caption}
\usepackage{enumitem}

\captionsetup{
  figurename=ภาพที่,
  tablename=ตารางที่
}

%-----------------------------------------------%
\setlength{\parindent}{\firstindentlength}
\setlist[enumerate]{
label=\thesection.\arabic*,
leftmargin=\parindent,
labelsep=1em,
itemindent=\parindent,
}
\makeatother
\setlist[itemize]{
leftmargin=\parindent,
labelsep=1em,
itemindent=\parindent,
}
\setlist[description]{
leftmargin=\parindent,
labelsep=1em,
itemindent=0.5\parindent,
}
%-----------------------------------------------%
% \setlist[itemize]{%
%   label=\textbullet
% }
% \setlist[enumerate]{label*=\arabic*.}

% \setlist[itemize,1]{%
%   leftmargin=*,
%   labelindent=\parindent,
%   labelsep=1em,
%   listparindent=0pt
% }
% \setlist[itemize,2]{%
%   leftmargin=*,
%   labelindent=0em,
%   labelsep=1em,
%   listparindent=0pt
% }
% \setlist[itemize,3]{%
%   leftmargin=*,
%   labelindent=0em,
%   labelsep=1em,
%   listparindent=0pt
% }

% \setlist[enumerate,1]{%
%   label=\thesection.\arabic*.,
%   leftmargin=*,
%   labelindent=\parindent,
%   labelsep=1em,
%   listparindent=0pt
% }
% \setlist[enumerate,2]{%
%   label=\thesection.\arabic{enumi}.\arabic*.,
%   leftmargin=*,
%   labelindent=0em,
%   labelsep=1em,
%   listparindent=0pt
% }
% \setlist[enumerate,3]{%
%   label=\thesection.\arabic{enumi}.\arabic{enumii}.\arabic*.,
%   leftmargin=*,
%   labelindent=0em,
%   labelsep=1em,
%   listparindent=0pt
% }

% \setcounter{secnumdepth}{3}
% \renewcommand\thesection{\arabic{section}}
% \renewcommand\thesubsection{\thesection.\arabic{subsection}}
% \renewcommand\thesubsubsection{\thesubsection.\arabic{subsubsection}}

% \titleformat{\section}
%   {\bfseries\large\upshape}
%   {\thesection}
%   {0.45\firstindentlength}
%   {}

% \titleformat{name=\section,numberless}
%   {\bfseries\large\upshape}
%   {}
%   {0pt}
%   {}

% \titleformat{\subsection}
%   {\bfseries\large\upshape}
%   {\thesubsection}
%   {0.45\firstindentlength}
%   {}

% \titleformat{name=\subsection,numberless}
%   {\bfseries\large\upshape}
%   {}
%   {0pt}
%   {}

% \titleformat{\subsubsection}
%   {\bfseries\large\upshape}
%   {\thesubsubsection}
%   {0.45\firstindentlength}
%   {}

% \titleformat{name=\subsubsection,numberless}
%   {\bfseries\large\upshape}
%   {}
%   {0pt}
%   {}

% \titlespacing{\section}
%   {0pt}
%   {\baselineskip}
%   {0\baselineskip}

% \titlespacing{\subsection}
%   {\firstindentlength}
%   {\baselineskip}
%   {0\baselineskip}

% \titlespacing{\subsubsection}
%   {\secondindentlength}
%   {\baselineskip}
%   {0\baselineskip}

% \pretocmd{\section}{\setlength{\parindent}{\firstindentlength}}{}{}
% \pretocmd{\subsection}{\setlength{\parindent}{\secondindentlength}}{}{}
% \pretocmd{\subsubsection}{\setlength{\parindent}{\thirdindentlength}}{}{}